% Drawing 43/04-12 -- folder 43, batch 04, archive page 65 (PDF page 66).
% Transcription: au centre, un croquis : un carré en pointillé dont les sommets sont des points, parcouru de flèches (une flèche descendante sur chacun des deux côtés verticaux, deux flèches vers la gauche sur le côté du bas), avec un arc à gauche et un arc au-dessous.
% Kind: dotted square with arrows. Method 2 (schematic redraw from the facsimile). Status: ok.
% Notes: Ink. Dots and short dashes mark out a square with a second column at its right; an arrow points down on the left side and one on the middle vertical, two arrows point left along the bottom; a bracket-like arc stands at the left and an arc under the bottom side. Still differs: the square is only suggested on the page by dots and dashes, and their number and places are approximate; which points the arrows join is a reading. Redrawn 2026-11-12 (issue 128): to check again. 2026-11-13, Claude Opus 5.5, overlay check (issue 147): redrawn on millimetre grids of the 228 ppi scan (rendered at 350 dpi in colour); not compiled (no TeX engine in the session), the TikZ was rasterised by a script and laid over the crop; every mark lies on its ink to within about 1 unit (the earlier ones were 2 to 10 units off). Drawn as on the page: the left column is a dot, a short thin stroke ending in a blot, a dot, an arrow down and a short dash; at the top of the second column a small bracket, a thin stroke and a dot, then an arrow down to a blot; along the bottom two heavy arrows to the left, each ending at a mark (the blot under the second column, a dot at the right end); the arc at the left (a bracket-like curve beside the first column) and the arc under the bottom row. So the dotted square has a third column at the right only through the right arrow and its end dot. Which marks are the vertices is the transcription's reading; the redraw gives the marks where they are.
% Coordinates: 1 unit = 1 mm = 1/100 of the crop width, y downwards as on the crop (bbox in page fractions: [0.33, 0.46, 0.54, 0.585]).
% Redrawn by Claude (claude-fable-5-2), issue #128, 2026-11-12; not checked by a human; not compiled (no TeX in the session).
% Rights: reproduces part of the work -- behind RENDER_TRANSCRIPTS, never published by itself.
\documentclass[tikz,border=3pt]{standalone}
\usepackage{amsmath,amssymb}
\usetikzlibrary{patterns,arrows.meta}
\tikzset{pencil/.style={black!45}, struck/.style={black!35}, mpt/.style={circle,draw,inner sep=1.3pt}, pt/.style={circle,fill,inner sep=0.9pt}, opt/.style={circle,draw,fill=white,inner sep=1.1pt}, Rd/.style={red!75!black}, Gn/.style={green!45!black}, Br/.style={brown!70!black}}
\begin{document}
\begin{tikzpicture}[x=1mm,y=-1mm,line cap=round,line join=round,font=\small]
% left column: dot, a short thin stroke ending in a blot, dot, arrow down, dash
\fill (25.6,13.4) circle[radius=0.6];
\draw (24.4,21.7) -- (25.2,28.9);
\fill (25.4,29.3) circle[radius=0.5];
\fill (25.9,38.3) circle[radius=0.6];
\draw[thick,-{Stealth[length=1.6mm]}] (26.4,45) -- (26.9,52.3);
\draw[thick] (26,60.8) -- (27.8,61.2);
% arc at the left
\draw[thick] plot[smooth] coordinates {(20,24.4) (18,28.5) (17.2,33) (17.4,38) (18.4,42.5) (20.5,47.5) (23,51)};
% top of the right column: a small bracket, a thin stroke, a dot
\draw[thick] plot[smooth] coordinates {(43.1,38.8) (41.9,39.9) (42.8,41.4)};
\draw (45,40) -- (53.3,39.4);
\fill (58,38.4) circle[radius=0.6];
% arrow down in the middle column
\draw[thick,-{Stealth[length=1.6mm]}] (55.8,44) -- (56.1,53.6);
% bottom row: two heavy arrows to the left between marks
\draw[very thick,-{Stealth[length=2.4mm]}] (48.3,61.4) -- (38,62.5);
\fill (55.2,60.6) circle[radius=0.6];
\draw[very thick,-{Stealth[length=2.4mm]}] (81.7,60.6) -- (68.3,62.5);
\fill (85,60.8) circle[radius=0.6];
% arc below
\draw[thick] plot[smooth] coordinates {(50.3,73.6) (56,75) (63,75.1) (70,73.3) (74.5,71.6) (76.4,70.6)};
\end{tikzpicture}
\end{document}
